\documentclass[12pt]{article}

% Layout.
\usepackage[top=1.2in, bottom=0.75in, left=1in, right=1in, headheight=1.0in, headsep=0pt]{geometry}

% Fonts.
\usepackage{mathptmx}
\usepackage[scaled=0.86]{helvet}
\renewcommand{\emph}[1]{\textsf{\textbf{#1}}}

% TiKZ.
\usepackage{tikz, pgfplots}
\usetikzlibrary{calc}
%\pgfplotsset{my style/.append style={axis x line=middle, axis y line=middle, xlabel={$x$}, ylabel={$y$}}}
%\pgfplotsset{compat=1.16}

% Misc packages.
\usepackage{amsmath,amssymb,latexsym}
\usepackage{graphicx}
\usepackage{array}
\usepackage{xcolor}
\usepackage{multicol}

% Commands to set various header/footer components.
\makeatletter
\def\doctitle#1{\gdef\@doctitle{#1}}
\doctitle{Use {\tt\textbackslash doctitle\{MY LABEL\}}.}
\def\docdate#1{\gdef\@docdate{#1}}
\docdate{Use {\tt\textbackslash docdate\{MY DATE\}}.}
\def\doccourse#1{\gdef\@doccourse{#1}}
\let\@doccourse\@empty
\def\docscoring#1{\gdef\@docscoring{#1}}
\let\@docscoring\@empty
\def\docversion#1{\gdef\@docversion{#1}}
\let\@docversion\@empty
\makeatother

% Headers and footers layout.
\makeatletter
\usepackage{fancyhdr}
\pagestyle{fancy}
\fancyhf{} % Clears all headers/footers.
\lhead{\emph{\@doctitle\hfill\@docdate}
\ifnum \value{page} > 1\relax\else\\
\emph{Name: \rule{3.5in}{1pt}\hfill \@docscoring}
\fi}

\rfoot{\emph{\@docversion}}
\lfoot{\emph{\@doccourse}}
\cfoot{\emph{\thepage}}
\renewcommand{\headrulewidth}{0pt}%
\makeatother

% Paragraph spacing
\parindent 0pt
\parskip 6pt plus 1pt

% A problem is a section-like command. Use \problem{5} for a problem worth 5 points.
\newcounter{probcount}
\newcounter{subprobcount}
\setcounter{probcount}{0}
\newcommand{\problem}[1]{%
\par
\addvspace{4pt}%
\setcounter{subprobcount}{0}%
\stepcounter{probcount}%
\makebox[0pt][r]{\emph{\arabic{probcount}.}\hskip1ex}\emph{[#1 points]}\hskip1ex}
\newcommand{\thesubproblem}{\emph{\alph{subprobcount}.}}

% Subproblems are an enumerate-like environment with a consistent
% numbering scheme. 
% Use \begin{subproblems}\item...\item...\end{subproblems}
\newenvironment{subproblems}{%
\begin{enumerate}%
\setcounter{enumi}{\value{subprobcount}}%
\renewcommand{\theenumi}{\emph{\alph{enumi}}}}%
{\setcounter{subprobcount}{\value{enumi}}\end{enumerate}}

% Blanks for answers in normal and math mode.
\newcommand{\blank}[1]{\rule{#1}{0.75pt}}
\newcommand{\mblank}[1]{\underline{\hspace{#1}}}
\def\emptybox(#1,#2){\framebox{\parbox[c][#2]{#1}{\rule{0pt}{0pt}}}}

% Misc.
\renewcommand{\d}{\displaystyle}
\newcommand{\ds}{\displaystyle}


\doctitle{Math 252: Quiz 1}
\docdate{Spring 2026}
\doccourse{}
\docversion{}
\docscoring{\fbox{{\LARGE \strut}\blank{0.8in} / 24}}

\begin{document}
24 points possible; each part is worth 2 points.  No aids (book, notes, calculator, phone, etc.) are permitted.  Show all work and use proper notation for full credit.  

\problem{12}  Compute the derivatives of the following functions.


\begin{subproblems}
%
\item   $\ds f(x) =x^2e^{x^2+1}+x+1$
\vfill

%basic rules they were told in advance they needed to know
\item  $\ds f(x)=\tan(2x)+\sin^2(x)\cos(x)$
\vfill



%
\item  $\ds f(x) = \frac{3x^2-1}{\ln(x^2+1)}$

\vfill

\newpage
\quad \\
% 
\item   $\ds f(x) =\sqrt{x+\sin^2(x^2)}$
\vfill


%
\item  $\ds f(x) = e^{1/x}-\frac{x^2}{2}+\frac{2}{\sqrt x}$
\vfill


% 
\item  $\ds f(x) = \ln(a+b^{2x})$ where $a$ and $b$ are constants
\vfill
\end{subproblems}


\newpage
\problem{12}  Compute the following antiderivatives (indefinite integrals) and definite integrals.  Remember that antiderivatives need a ``$+C$''.


\begin{subproblems}


% 1.2 #80
\item $\ds \int_0^\pi 4e^x+\cos(x) \; dx$
\vfill


% like 1.3 #172
\item $\ds \int_0^3 (x+2)(x-3)\; dx$
\vfill


%
\item $\ds \int 2x^3\sqrt{x^2+1}\; dx$
\vfill


\newpage
\quad


%% like practice quiz 1 #f
\item $\ds \int_0^{\pi/3} \sec^2(x)\tan(x) \; dx$
\vfill


%% HW 1.7 # 424 
\item $\ds \int \frac{e^x}{1+e^{(2x)}} \; dx$
\vfill

%% HW 1.6 # 348
\item $\ds \int \tan(2x) \; dx$

\vfill


\end{subproblems}
\end{document}